\subsection{Sub-Phase 3.3: Geometrically Nonlinear Dynamics and Time Integration}
\label{subsec:dynamics_time}

To address dynamic loading scenarios such as structural vibration, transient shocks, and time-dependent cyclic forces, the physics engine was extended from a static balance state to a fully dynamic continuum representation. The semidiscrete governing equation of motion in physical space is expressed as:
\begin{equation}
        \mat{M} \ddot{\vect{u}}(t) + \mat{C} \dot{\vect{u}}(t) + \vect{F}_{\text{int}}(\vect{u}(t)) = \vect{F}_{\text{ext}}(t)
\end{equation}
where $\mat{M} \in \mathbb{R}^{n \times n}$ represents the consistent global mass matrix, $\mat{C} \in \mathbb{R}^{n \times n}$ represents the viscous damping matrix, $\vect{F}_{\text{int}}(\vect{u})$ is the geometrically nonlinear internal force vector, and $\vect{F}_{\text{ext}}(t)$ represents the time-dependent external load vector.

The global mass matrix $\mat{M}$ is assembled element-wise from consistent element mass matrices $\mat{M}_e$ computed via Gauss-Legendre integration:
\begin{equation}
        \mat{M}_{e, ab} = \int_{-1}^{1} \int_{-1}^{1} \rho R_a(\xi, \eta) R_b(\xi, \eta) |J(\xi, \eta)| \, d\xi \, d\eta
\end{equation}
where $\rho$ is the material mass density, and $R_a$ represents the bivariate NURBS shape functions. Damping is modeled using a standard Rayleigh formulation:
\begin{equation}
        \mat{C} = \alpha_M \mat{M} + \beta_K \mat{K}_0
\end{equation}
where $\alpha_M$ and $\beta_K$ are user-specified mass and stiffness proportional damping factors, and $\mat{K}_0$ is the baseline linear stiffness matrix.

\subsubsection{Reduced-Order Dynamic Equations}
By projecting the displacement state onto the low-dimensional orthogonal POD basis ($\vect{u}(t) \approx \mat{V}_r \vect{q}_r(t)$), the dynamic equations of motion pull back to the reduced state-space manifold of dimension $r \ll n$:
\begin{equation}
        \mat{M}_r \ddot{\vect{q}}_r(t) + \mat{C}_r \dot{\vect{q}}_r(t) + \mat{V}_r^T \vect{F}_{\text{int}}(\mat{V}_r \vect{q}_r(t)) = \vect{F}_{r,\text{ext}}(t)
\end{equation}
where the pre-projected reduced mass and damping matrices are invariant:
\begin{equation}
        \mat{M}_r = \mat{V}_r^T \mat{M} \mat{V}_r, \quad \mat{C}_r = \mat{V}_r^T \mat{C} \mat{V}_r, \quad \vect{F}_{r,\text{ext}}(t) = \mat{V}_r^T \vect{F}_{\text{ext}}(t)
\end{equation}

\subsubsection{\texorpdfstring{Newmark-$\beta$}{Newmark-beta} Implicit Time-Stepping and Tangent System}
To resolve the second-order nonlinear differential system in time, an implicit Newmark-$\beta$ integration scheme is implemented. The displacement and velocity vectors at step $n+1$ (corresponding to time $t_{n+1} = t_n + \Delta t$) are approximated as:
\begin{align}
        \vect{u}_{n+1} &= \vect{u}_n + \Delta t \dot{\vect{u}}_n + \frac{\Delta t^2}{2} \left[ (1 - 2\beta)\ddot{\vect{u}}_n + 2\beta \ddot{\vect{u}}_{n+1} \right] \\
        \dot{\vect{u}}_{n+1} &= \dot{\vect{u}}_n + \Delta t \left[ (1 - \gamma)\ddot{\vect{u}}_n + \gamma \ddot{\vect{u}}_{n+1} \right]
\end{align}
where $\beta \in [0, 0.5]$ and $\gamma \in [0, 1.0]$ are the Newmark integration parameters (set to $\beta=0.25$ and $\gamma=0.5$ for unconditional stability and zero numerical dissipation). 

At each time step $n+1$, the dynamic residual vector $\vect{R}_{\text{dyn}}$ must vanish:
\begin{equation}
        \vect{R}_{\text{dyn}}(\vect{u}_{n+1}) = \vect{F}_{\text{ext}}^{n+1} - \mat{M} \ddot{\vect{u}}_{n+1} - \mat{C} \dot{\vect{u}}_{n+1} - \vect{F}_{\text{int}}(\vect{u}_{n+1}) = \vect{0}
\end{equation}
Inserting the Newmark relations and linearizing around iteration $k$ yields the dynamic tangent system:
\begin{equation}
        \mat{K}_T^{\text{dyn}}(\vect{u}_{n+1}^k) \Delta \vect{u} = \vect{R}_{\text{dyn}}(\vect{u}_{n+1}^k)
\end{equation}
where the dynamic tangent stiffness matrix $\mat{K}_T^{\text{dyn}}$ incorporates the inertial and damping contributions:
\begin{equation}
        \mat{K}_T^{\text{dyn}} = \mat{K}_T(\vect{u}_{n+1}^k) + \frac{1}{\beta \Delta t^2} \mat{M} + \frac{\gamma}{\beta \Delta t} \mat{C}
\end{equation}
The equivalent projection of this linearized step to the reduced manifold reads:
\begin{equation}
        \left( \mat{K}_{r,T} + \frac{1}{\beta \Delta t^2} \mat{M}_r + \frac{\gamma}{\beta \Delta t} \mat{C}_r \right) \Delta \vect{q}_r = \mat{V}_r^T \vect{R}_{\text{dyn}}(\mat{V}_r \vect{q}_{r,n+1}^k)
\end{equation}
This formulation preserves unconditional stability and allows the client-side simulator to update dynamic trajectories interactively.
